El proyecto \textit{m2r} es una herramienta educativa diseñada para la enseñanza de los principios fundamentales de la compilación de lenguajes de programación. Su objetivo principal es permitir al estudiante comprender el proceso de traducción de un lenguaje de alto nivel a un lenguaje intermedio ejecutable por una máquina virtual.

Este sistema se divide en dos componentes principales:
\begin{enumerate}
    \item \textbf{El Compilador:} Encargado de realizar el análisis léxico, sintáctico y semántico del código fuente, generando un código intermedio.
    \item \textbf{La Máquina Virtual (Intérprete m2r):} Encargada de ejecutar el código generado por el compilador, simulando un entorno de ejecución con memoria, registros y un ciclo de instrucción definido.
\end{enumerate}

Este documento aborda los fundamentos teóricos detrás del diseño de compiladores, el detalle técnico de la implementación de este proyecto, y una guía práctica para su uso y extensión.
